\documentclass[11pt,letterpaper]{article}

\usepackage[margin=1in]{geometry}
\usepackage[T1]{fontenc}
\usepackage[utf8]{inputenc}
\usepackage{XCharter}
\usepackage[expansion=alltext,protrusion=true]{microtype}
\usepackage{xcolor}
\usepackage{hyperref}

\definecolor{cvblue}{RGB}{18,84,178}
\hypersetup{
  colorlinks=true,
  urlcolor=black,
  pdftitle={Amrita - Statement of Purpose},
  pdfauthor={Amrita},
  pdfsubject={PhD in Educational Research, The University of Alabama}
}

\pagestyle{plain}
\linespread{1.25}
\setlength{\parindent}{0pt}
\setlength{\parskip}{9pt}
\widowpenalty=10000
\clubpenalty=10000

\begin{document}

\begin{center}
  {\Large\MakeUppercase{Amrita}}\\[5pt]
  {\small \href{mailto:hiamritaa@gmail.com}{hiamritaa@gmail.com}}\\[8pt]
  {\large Statement of Purpose}\\[3pt]
  {\small PhD in Educational Research (Qualitative Research), The University of Alabama}
\end{center}
\vspace{-4pt}
\noindent\textcolor{cvblue}{\rule{\linewidth}{0.5pt}}
\vspace{2pt}

I grew up in a defence family, which meant moving from one place to another since I was very young. I've stayed in villages with almost no resources and in five-star hotels with everything in them, and I believe a big part of how I understand the world comes from that. It is probably also why the question I keep coming back to is how people in very different circumstances live with the same machines.

Recently I put a version of that question into a survey. I got 156 responses across India on whether a machine failing at an important moment ever feels like more than a technical failure. In big cities, 46\% said it often or sometimes does. In smaller towns and semi-rural areas, it was between 70 and 74\%. I still can't say what that gap means. People outside the cities may simply relate to machines differently. Or the question meant something different to each group, or the Hindi scenarios I read aloud in interviews didn't match the written ones. With the training I have now, I can't tell these apart, and that is a big part of why I'm applying to the PhD in Educational Research at The University of Alabama.

I studied Fashion Communication, a communication design degree, at the National Institute of Fashion Technology (NIFT) and graduated in 2024 with a CGPA of 8.3 out of 10. Along the way I was ranked first in India in NIFT's selection for its dual-degree exchange with the Fashion Institute of Technology in New York, where I spent a year studying advertising and marketing communications. That year taught me how critical information gets shaped to suit an advertising goal. A shop can promise ``Buy a bike, get a free car'' when the small print says the car is only a lucky draw. E-commerce platforms and shopping ads do this all the time, and it inspired me to write my paper on deceptive dark patterns in Indian fashion e-commerce.

Design is where I found research. In 2022 our team spent several days in Raiyam, a village in Bihar, documenting Sikki, a grass-weaving craft, as part of our coursework. I interviewed three generations of one artisan family and several other artisans, recorded their processes and techniques, and then compared what the family told me with the craft's official Geographical Indication record. Sitting with people and then checking their account against an official one is still how I like to work.

Since graduating, I've balanced a full-time design job with research, and I have three papers accepted, at India HCI 2026, HWWE 2026 and FTSC 2027. For the FTSC paper, written with Prof.~Ragini Ranjana of NIFT Patna, I looked at how Indian e-commerce sites use heritage crafts in festival marketing. Not one of the 75 product pages I studied named the artisan who made the item. My latest study, now under review at \textit{Technology in Society}, looks at how people in India relate to their machines. It combines the survey with 21 interviews, and 96\% of the people I surveyed reported some form of care for their machines, such as blessing a new vehicle or their tools before first use. I went into the interviews expecting people to talk about machines as if they were alive. Most of them described that care as plain maintenance. Getting it wrong taught me more than the survey numbers did, and it showed me how much the design of an interview decides what a researcher gets to hear.

My goal is to understand how people across different socioeconomic contexts live with, trust and make sense of everyday machines and new technology, including AI tools. I see this as a question about learning. Someone blessing their tools at a festival, or a person with dyslexia trying a prompt-based image generator, is quietly working out how the system works and how far to trust it, usually on their own. Teaching children with the Robin Hood Army and being part of Teach For India's School Design Cohort showed me how much learning happens outside a classroom, and made me curious about how people learn to live with technology in everyday life.

What I don't have yet is methodological depth. I can run a survey and an interview study, but I can't check whether a question comes across the same way in different languages, or in a city and a village, and I've read too little qualitative methodology to design interviews that do more than collect answers. I also ran the machines study on my own, without any guidance, and I want to learn to do this kind of work properly. That is what I'm hoping the program will teach me.

The University of Alabama is where I'd like to do that. Dr.~Susan Nordstrom's work on how nonhuman objects, including technological ones, take part in qualitative research is the approach I most want to learn. Her argument that even the recorder on the table shapes an interview fits interviews like mine, where the machine was the subject of every conversation, and I would like to work with her as my advisor. I also want enough measurement training to answer the question I started with. Dr.~Stefanie Wind studies whether survey questions and rating scales work the same way across groups, which is what I couldn't check in my city and small-town comparison. Her methods are statistical and new to me, which is part of why I want to learn them here.

After the PhD I want to teach research methods and keep doing this work in India, with studies that hold up across languages and income groups instead of assuming a question means the same thing everywhere. I'm applying for an assistantship and for external scholarships. I'm used to fitting research around a full-time job, I adapt quickly to unfamiliar places, and as a communication designer I can turn findings into something people outside academia can understand. I could be a valuable asset to UA since I am eager to learn more and, my thirst for information and possibilities continues to grow.

\end{document}
